\SourceLeaf{62}glückliche Erfolge bei drei andern seitdem entdeckten neuen Planeten auf das vollkommenste weggeräumt.

Jene Methoden sind die ersten Grundlagen des Werks, dem diese Anzeige gewidmet ist: allein die fortgesetzte Beschäftigung mit diesem Gegenstande hat dem Verfasser Gelegenheit gegeben, so viele wiederholte Zusätze und Abänderungen und Vervollkommnungen an denselben anzubringen, dass von ihrer ursprünglichen Gestalt fast gar keine Spuren übrig geblieben sind. Der Verfasser hofft daher, dass die Astronomen, welche von Anfang an ihr lebhaftes Interesse an diesem Untersuchungen und den Wunsch nach einer baldigen Bekanntmachung geäussert haben, mit der verspäteten Erscheinung nicht unzufrieden zu sein Ursache haben werden. Indem wir das Urtheil, in wie fern diese Hoffnung gegründet ist, den competenten Lesern des Werks selbst überlassen, begnügen wir uns hier, nur den Plan und Inhalt desselben in gedrängter Kürze darzulegen.

Jenes mehrerwähnte Problem war freilich die Veranlassung und der Hauptzweck des vorliegenden Werks: allein in der Ausführung konnte dasselbe doch nur ein Theil, ja nur der kleinere Theil desselben werden. Man wird leicht vermuthen, dass eine adäquate Auflösung desselben sich nicht geben liess, ohne eine Menge bisher noch unentwickelter Wahrheiten vorauszusetzen, die sich auf die Bewegung der Himmelskörper in Kegelschnitten beziehen. Diese neuen, auch an sich selbst schon sehr interessanten, Relationen waren so zahlreich, und standen unter sich in so inniger Verbindung, dass sie nothwendig im Zusammenhange vorgetragen werden mussten, daher ihnen die ganze erste Abtheilung des Werks gewidmet ist. Der Plan, wonach sie geordnet sind, ist sehr einfach: diese erste Abtheilung zerfällt in vier Abschnitte, wovon der erste und zweite die Relationen enthalten, die sich auf einen einzigen Ort des Himmelskörpers entweder in seiner Bahn, oder im Raume, beziehen; während der dritte und vierte Abschnitt die Relationen zwischen zweien oder mehreren Oertern entwickeln. Dass hier durchgehends neben den neuen Relationen auch die vornehmsten schon bekannten mit aufgeführt werden mussten, liess sich nicht wohl vermeiden: man wird indess letztere meistens in einer dem Verfasser eigenthümlichen Form entwickelt finden. Der beschränkte Raum dieser Blätter erlaubt es nicht, von allen hier abgehandelten Gegenständen auch nur eine Übersicht zu geben; wir schränken uns darauf ein, nur ein paar von denjenigen Untersuchungen zu berühren, die von etwas grösserem Umfange sind. Dahin gehört im ersten Abschnitte eine Digres%
\SourceLeaf{63}sion über den Grad der Genauigkeit, den man im numerischen Calcul bei dem Gebrauche von Tafeln aller Art zu erreichen im Stande ist; ferner in demselben Abschnitte eine neue Methode, die Bewegung in solchen Ellipsen und Hyperbeln, welche sich der Parabel nähern, auf die Bewegung in der letztern zu reduciren; im zweiten Abschnitte ein Verfahren, die geocentrischen Örter der Himmelskörper mit Hülfe dreier Coordinaten zu bestimmen, worüber bekanntlich der Verfasser schon früher eine Abhandlung geliefert hatte, welches aber hier noch eine wesentliche Vervollkommnung erhält; eben daselbst die Bestimmung der Differential-Änderungen des geocentrischen Orts durch die Differential-Änderungen der einzelnen Elemente. Mit vorzüglicher Ausführlichkeit ist im dritten Abschnitte das äusserst wichtige Problem abgehandelt, aus zweien Örtern eines Himmelskörpers in seiner Bahn die Elemente zu bestimmen, wo besonders die zweite Auflösung zur Entwickelung von einer grossen Anzahl neuer Relationen Gelegenheit gegeben hat. Um die practische Brauchbarkeit der vorgetragenen Methoden desto besser ins Licht zu setzen, und das Studium derselben auch den weniger Geübten mehr zu erleichtern, sind die wichtigern Lehren durchgehends mit Beispielen erläutert, die an wirklichen Fällen gewählt sind und grössten Theils wieder unter einander im Zusammenhange stehen. Die Bewegung in der Hyperbel, die doch auch bei Cometen wirklich vorkommen kann, ist schon des analytischen Interesses wegen beinahe mit derselben Ausführlichkeit abgehandelt, wie die Bewegung in der Ellipse und Parabel.

Die in der ersten Abtheilung vorgetragenen Untersuchungen sind nun gleichsam der Stoff, aus welchem in der andern Abtheilung die Auflösung der grossen Aufgabe, aus der Zergliederung der geocentrischen Erscheinungen die Bahn des Himmelskörpers selbst zu bestimmen, zusammengesetzt wird. Eigentlich beruht dieses Geschäft wieder auf mehreren unter sich sehr verschiedenen Forderungen. Eine andere Arbeit ist nöthig, um zum ersten Male die noch ganz unbekannte Bahn eines Himmelskörpers näherungsweise zu bestimmen; eine andere, um die schon näherungsweise bekannten Elemente nach einer längern, vielleicht viele Jahre umfassenden, Reihe von Beobachtungen auszufeilen. Die erstere Aufgabe erfordert grössere Kunst, die andere grössere Arbeit. Zu jener wird man nicht mehr Beobachtungen anwenden, als eben nöthig sind, also, allgemein zu reden, drei, und das Problem, aus dreien vollständigen Beobachtungen die Bahn eines Himmelskörpers zu bestimmen, gewisser Maassen das Wich%
\SourceLeaf{64}%
tigste des ganzen Werks, macht den Inhalt des ersten Abschnitts der zweiten
Abtheilung aus. Die Auflösung selbst verträgt hier keinen Auszug; als eine
Probe von der Allgemeinheit ihrer Anwendbarkeit führen wir nur an, dass sie
in den drei erläuternden Beispielen mit gleichem Erfolge auf Beobachtungen,
die 22 Tage, auf solche, die 71 Tage, und auf solche, die 260 Tage von einander
abstehen, angewandt ist. Etwas Eigenthümliches ist es, dass man nach
Gefallen die kleinen Modificationen, welche von der Parallaxe und Aberration
herrühren, \textit{sogleich} ohne bedeutende Erschwerung der Arbeit mit in Betrachtung
ziehen kann, und also ihretwegen nicht nöthig hat, eine doppelte Rechnung zu
führen.

Unter gewissen Umständen, besonders wenn die Ebene der Bahn wenig
gegen die Ecliptik geneigt ist, würde es nicht zweckmässig sein, die Bestimmung
der Elemente auf drei vollständige Beobachtungen zu gründen: man muss
vier Beobachtungen zum Grunde legen, worunter aber nur zwei vollständige zu
sein brauchen. So ergibt sich eine neue Aufgabe, welcher der zweite Abschnitt
gewidmet ist. Ein Beispiel, von der Vesta hergenommen, wo die äussersten Beobachtungen
162 Tage von einander entfernt liegen, erläutert diese Methode.
Auch hier kann man, wenn man will, \textit{sogleich} auf Parallaxe und Aberration
Rücksicht nehmen.

Zu der schärfern Ausfeilung der Elemente eines Himmelskörpers hat man,
gerade umgekehrt, nicht die möglich kleinste Zahl von Beobachtungen, sondern
so viele, als nur zu Gebote stehen, anzuwenden. Wie man sich dabei zu verhalten
habe, lehrt der dritte Abschnitt. Hier war der Ort, die Haupt-Momente
von einer für jede Anwendung der Mathematik auf die Körperwelt höchst wichtigen
Frage zu entwickeln, wie Beobachtungen und Messungen, die bei der Unvollkommenheit
unsrer Sinne und Werkzeuge unvermeidlich immer mit Fehlern,
wenn auch noch so geringen, behaftet sind, am zweckmässigsten zur Festsetzung
von Resultaten zu combiniren sind. Die Grundsätze, welche hier ausgeführt
werden, und welche von dem Verfasser schon seit 14 Jahren angewandt, und von
demselben schon vor geraumer Zeit mehreren seiner astronomischen Freunde mitgetheilt
waren, führen zu derjenigen Methode, welche auch \textsc{Legendre} in seinem
Werke: \textit{Nouvelles méthodes pour la détermination des orbites des comètes}, vor einigen
Jahren unter dem Namen \textit{Méthode des moindres carrés} aufgestellt hat: die
Begründung der Methode, welche von dem Verfasser gegeben wird, ist diesem
\SourceLeaf{65}%
ganz eigenthümlich. Eine weitere Ausführung hat man von demselben in der
Folge zu erwarten.

Um eine \textit{bleibende} Übereinstimmung der Elemente mit den Bewegungen eines
Planeten zu erreichen, darf man sich nicht auf eine rein elliptische Bewegung
einschränken: man muss dazu nothwendig die von den andern Planeten
hervorgebrachten Störungen mit in Betrachtung ziehen. Die Berechnung der Störungen
selbst lag natürlich ausser dem Plane eines Werks, welches nur der Bewegung
nach den \textsc{Kepler}schen Gesetzen gewidmet ist: allein von den nöthigen
Operationen, um eine Planetenbahn mit Rücksicht auf die schon berechneten
Störungen zu verbessern, und so gleichsam die lezte Hand an dieselbe zu legen,
sind im vierten und letzten Abschnitte wenigstens die Haupt-Momente angedeutet.

Wir schliessen diese Anzeige mit dem Wunsche, dass dieses Werk, die
Frucht einer siebenjährigen Beschäftigung mit diesen Untersuchungen, dazu beitragen
möge, die Freunde dieses wichtigsten und schönsten Theils der theorischen
Astronomie zu vermehren, und den Astronomen selbst ihre bei dem steten
Wachsthume der Wissenschaft an Umfang zunehmenden Arbeiten zu erleichtern.
\par\bigskip\centerline{\rule{.24\textwidth}{.3pt}}
